\section{Conclusions}

In this work, we introduced \ours{}, which enhances \sam{} to accept and propagate a \textit{visual} prompt across frames, and showed that its memory attention mechanism can be re-focused towards semantic correspondences by restructuring the feature space. We addressed three fundamental inquiries: (1) we demonstrated that semantic information can be extracted from \sam{} features through lightweight bottleneck transformations, answering our initial research question,  (2) we showed that this can be achieved while keeping \sam{} frozen, thereby preserving its segmentation capabilities, and finally (3) we experimentally verified that the learned semantic mapping generalizes robustly to novel classes, with SOTA performance on strict few-shot benchmarks and against generalist models. Beyond technical contributions, \ours{} offers practical advantages: it supports diverse prompts for annotation-efficient applications, and offers a 3x speed increase. Finally, our results suggest that foundation models like \sam{} may contain richer task-adaptable knowledge than their original objectives imply, a direction worthy of future exploration.


\newpage
